%======================================================================
\section{Desempeño con 50 realizaciones y las métricas de la tesis}\label{sec:desempeno}
%======================================================================
\paragraph{Por qué hace falta.} Las secciones anteriores usan entre 1 y 20 realizaciones. Bastan
para comparaciones pareadas con intervalos estrechos, pero un paper necesita cifras con
intervalos de confianza honestos, medidas con las mismas métricas con que la tesis evaluó sus
controladores.

\paragraph{La campaña F6.} Se corrieron 5850 lazos cerrados en la máquina Olab: 24 procesos
en paralelo, 175~min y ninguna corrida fallida (\code{julia/f6\_campaign.jl},
\code{julia/F6Tasks.jl}). Las campañas \code{nom} y \code{ext} de esta sección usan 50
realizaciones de la familia nominal (semillas impares 1--99, 121 consultas) y de la familia
extrema (semillas 101--150, 481 consultas). Las campañas de riesgo se describen en la
parte~VI. Todas las comparaciones usan números aleatorios comunes, y los intervalos son la
media $\pm1.96$ errores estándar. Tablas completas en \code{results/f6\_paper/summary.md}
(\code{python/f6\_summary.py}).

\paragraph{Las métricas de la tesis} (Tabla~10 del paper de control; \code{thesis\_metrics} en
\code{julia/ClosedLoop.jl}). Con $e=y_{\rm ref}-\mathrm{PCT}$:
\begin{itemize}
  \item \textbf{IAE}: integral de $|e|$ en el tiempo; resume el error de seguimiento.
  \item \textbf{Tiempo de asentamiento}: instante desde el cual $|e|\le0.30$~K para siempre
        ($\infty$ si nunca ocurre). Se reporta la mediana y la fracción de realizaciones que
        se asientan.
  \item \textbf{Sobrepaso mínimo}: $\min(\mathrm{PCT}-y_{\rm ref})$; negativo si la PCT baja
        de la consigna (sobre-enfriamiento).
  \item \textbf{Error estacionario}: media de $e$ en el último 20\,\% de la corrida.
  \item \textbf{Variación total (VT)}: $\sum|\Delta\dot m|$ del caudal mandado; mide cuánto se
        mueve el actuador.
  \item \textbf{DSM mínimo}: margen mínimo a 1244~K, el límite de la tesis.
\end{itemize}

\begin{table}[H]\centering\scriptsize\setlength{\tabcolsep}{3pt}
\begin{tabularx}{\linewidth}{@{}Y r r r r r r@{}}\toprule
Controlador & IAE [K\,s] & Asienta (fracción) & Sobrepaso mín.\ [K] & Err.\ estac.\ [K] & VT [kg/s] & DSM mín.\ [K]\\\midrule
\multicolumn{7}{@{}l}{\textit{Familia nominal, bomba archivada (12.9~kg/s)}}\\
N1 = N2 = N3 (+sesgo) & $14.7\pm1.6$ & 13~s (100\,\%) & $-0.02$ & 0.01 & 30 & 62.9\\
PID de la tesis & $92.9\pm0.4$ & nunca & $-2.73$ & 1.99 & 15\,340 & 62.9\\
PI del DET & $221\pm13$ & nunca & $+1.84$ & $-5.18$ & 0 & 59.2\\
\addlinespace
\multicolumn{7}{@{}l}{\textit{Familia nominal, 40~kg/s, rampa 5~kg/s por paso}}\\
N1 = N2 = N3 (+sesgo) & $7.6\pm0.5$ & 5.5~s (100\,\%) & $-0.11$ & 0.01 & 87 & 62.9\\
PID de la tesis & $208\pm1$ & nunca & $-5.20$ & 3.11 & 15\,060 & 62.9\\
PI del DET & $221\pm13$ & nunca & $+1.84$ & $-5.18$ & 0 & 59.2\\
\addlinespace
\multicolumn{7}{@{}l}{\textit{Familia extrema, actuador físico (ninguna corrida llega a 1220~K)}}\\
N2 = N3 (+sesgo) & $919\pm280$ & 85~s (94\,\%) & $-6.1$ & 0.79 & 468 & 50.0\\
N1 (+sesgo) & $637\pm110$ & 86~s (100\,\%) & $-2.1$ & $-0.05$ & 418 & 50.0\\
N2 sin sesgo & $835\pm100$ & 100~s (14\,\%) & $+0.74$ & $-0.76$ & 324 & 49.1\\
PID de la tesis & $3654\pm800$ & 238~s (2\,\%) & $-33.3$ & 12.6 & 9576 & 45.3\\
PI del DET & $5755\pm49$ & nunca & $-40.7$ & 24.5 & 972 & 35.2\\\bottomrule
\end{tabularx}
\caption{Métricas de la tesis en 50 realizaciones por familia (media $\pm$ IC al 95\,\%). Con
la bomba archivada en la familia extrema, N2 llega a 1220~K en 22 de 50 realizaciones porque la
capacidad no alcanza (sección~\ref{sec:extremos}). La VT del PID es grande porque su caudal
oscila entre el mínimo y el máximo de la bomba.}
\label{tab:desempeno}
\end{table}

\paragraph{Cómo leer la tabla.}
\begin{enumerate}
  \item \textbf{En el escenario de la tesis (familia nominal), N2 gana en todas las métricas.}
        Se asienta en todas las realizaciones, en 5.5--13~s, con error estacionario de 0.01~K.
        El PID de la tesis nunca se asienta en la banda de 0.30~K: sobre-enfría (sobrepaso
        de $-2.7$ a $-5.2$~K) y su caudal oscila (VT de unos 15\,000~kg/s). El PI del DET
        trabaja con su propia consigna de 1144~K y deja la PCT 5~K sobre $y_{\rm ref}$.
  \item \textbf{En la familia extrema, N1, N2 y N3 dan la misma PCT pico en las 50
        realizaciones}, y la diferencia mediana de IAE entre N1 y N2 es de $-0.1$~K\,s. La
        media difiere ($-282\pm170$~K\,s a favor de N1) por unas pocas realizaciones en las que,
        después de una falla larga con $\eta$ y sesgo de pronóstico desconocidos, las
        trayectorias se separan. Optimizar mejor el costo P0 no garantiza un IAE menor cuando
        la planta tiene perturbaciones que el modelo no conoce. La corrección de sesgo hace que
        N2 se asiente (94\,\% frente a 14\,\% sin ella).
  \item \textbf{El PID y el PI son los que más arriesgan.} Su DSM mínimo es de 45 y 35~K,
        frente a 50~K de N2, y su IAE es de 4 a 6 veces el de N2.
\end{enumerate}

\takeaway{\textbf{Respuesta a Q6 y Q8.} Con 50 realizaciones por familia, las conclusiones
se sostienen. N1, N2 y N3 dan la misma PCT pico, y N2 supera al PID y al PI de la tesis en
todas sus métricas: se asienta siempre en el escenario de la tesis, con un IAE de 6 a 28 veces
menor que el del PID. El margen DSM de la tesis no cambia la trayectoria de N2
(sección~\ref{sec:dsm-lambda}).}

%======================================================================
\section{¿Los solvers resuelven a optimalidad en lazo cerrado?}\label{sec:auditoria}
%======================================================================
La certificación de la sección~\ref{sec:certificacion} vale para dos consultas. En lazo
cerrado hay miles, y en cada una el solver puede terminar antes de tiempo o en un óptimo local
distinto. La auditoría responde a esa pregunta en dos niveles: primero, qué dicen los estados
de terminación de los solvers; después, cuán lejos queda cada decisión del mejor plan que se
puede encontrar.

\subsection{Estados de terminación}
El conteo de fallas del Monte Carlo trata como éxito la salida \emph{aceptable} de Ipopt
(\code{ALMOST\_LOCALLY\_SOLVED}). Para saber si eso esconde algo, a lo largo de un lazo
N3+sesgo (726 consultas: la realización archivada y dos del Monte Carlo, con los dos
actuadores) cada consulta se resolvió de cinco maneras, y todas las soluciones se evaluaron
con el mismo costo P0 sobre el ROM (\code{julia/f5a\_optimality\_audit.jl}).

\begin{table}[H]\centering\footnotesize\setlength{\tabcolsep}{4pt}
\begin{tabularx}{\linewidth}{@{}l Y r r@{}}\toprule
Solver & Estado de terminación (726 consultas) & Brecha [K$^2$] (mediana / máx.) & $|\Delta\mdot_0|$ máx.\ [kg/s]\\\midrule
N1 (QP) & óptimo: 726 & $10^{-6}$ / 0.0097 & 0.11\\
N2 (QP) & óptimo: 726 & $5\cdot10^{-12}$ / 0.0022 & 0.01\\
N3 en el lazo & óptimo local: 160; aceptable: 566 & $2\cdot10^{-14}$ / $3\cdot10^{-8}$ & ---\\
N3 sin salida aceptable & igual, más 7 por límite de iteraciones & idéntica & $2\cdot10^{-7}$\\
N3 con arranque en frío & igual & 0 / $4\cdot10^{-8}$ & 0.0015\\\bottomrule
\end{tabularx}
\caption{Estados de terminación. La brecha de costo se mide contra el mejor N3 de cada
consulta. Estados: óptimo = \code{OPTIMAL} (Clarabel); óptimo local = \code{LOCALLY\_SOLVED}
y aceptable = \code{ALMOST\_LOCALLY\_SOLVED} (Ipopt).}
\end{table}

\noindent N1 y N2 resuelven su QP convexo al óptimo global. Ese óptimo es el del modelo
linealizado; medido sobre P0, la pérdida es $\le0.01$~K$^2$ y sólo aparece en las primeras
consultas, donde el costo vale 15--46~K$^2$. N3 termina en la salida aceptable en el 78\,\% de
las consultas porque, en régimen, el costo óptimo es $\approx3\cdot10^{-12}$~K$^2$: el
seguimiento es perfecto y la tolerancia de $10^{-10}$ queda bajo el ruido numérico. Sin la
salida aceptable, o desde un arranque en frío, llega al mismo punto.

\subsection{Distancia al mejor plan, en 3800 consultas por lazo}
En la campaña F6, a lo largo de los lazos de N1, N2 y N3 de las tres familias, se auditó una
consulta de cada 20 (de cada 10 en la nominal): 3800 consultas por cada lazo, unas 11\,400 en
total. En cada una se
resolvieron N1, N2, N3 con arranque en caliente y pulido, y N3 desde un arranque en frío. Todos
los planes se evaluaron con el mismo costo P0 sobre el ROM \code{ode1be}, y $J_{\rm best}$ es el
mejor de ellos.

\paragraph{Cómo medir la distancia: una lección de normalización.} La medida natural sería la
brecha relativa $J/J_{\rm best}-1$, pero en el 53--85\,\% de las consultas la consigna se sigue
exactamente y $J_{\rm best}\approx10^{-12}$~K$^2$. Dividir por un número así infla cualquier
diferencia: la primera versión de la tabla mostraba brechas de $10^5$\,\% entre planes que
daban la misma trayectoria. Por eso el exceso se mide como un error de seguimiento equivalente,
en kelvin: $\Delta_{\rm RMS}=\sqrt{J/K}-\sqrt{J_{\rm best}/K}$ con $K=26$. La brecha relativa se
reporta sólo donde $J_{\rm best}\ge1$~K$^2$.

\begin{table}[H]\centering\footnotesize\setlength{\tabcolsep}{4pt}
\begin{tabularx}{\linewidth}{@{}l r r Y@{}}\toprule
Lazo (familia, actuador) & Consultas & $\Delta_{\rm RMS}$ máx.\ [K] & Brecha relativa donde $J_{\rm best}\ge1$ (mediana / máx.)\\\midrule
N3, todas las familias & 3800 & $2.5\cdot10^{-8}$ & $\le3\cdot10^{-9}$\\
N2, con fallas del DET (40~kg/s) & 1250 & $6.3\cdot10^{-5}$ & $2\cdot10^{-9}$ / $8\cdot10^{-5}$\\
N2, nominal (archivada / 40~kg/s) & 650 / 650 & $3.8\cdot10^{-6}$ / $1.0\cdot10^{-4}$ & $\le4\cdot10^{-4}$\\
N2, extrema (físico) & 1250 & $5.9\cdot10^{-4}$ & $5\cdot10^{-6}$ / $3\cdot10^{-4}$\\
N1, con fallas y nominal & 2550 & $\le6.8\cdot10^{-4}$ & $\le8\cdot10^{-4}$\\
N1, extrema (físico) & 1250 & 0.33 & $4\cdot10^{-3}$ / 0.063\\\bottomrule
\end{tabularx}
\caption{Distancia al mejor plan encontrado (\code{results/f6\_paper/*\_audit.csv}). N3 con
arranque en frío no encontró un plan mejor en ninguna consulta.}
\end{table}

\noindent N1 linealiza una sola vez, en el punto de la consulta. Cuando el plan óptimo hace
movimientos grandes, como en la familia extrema con el actuador físico, el modelo lineal se
aleja del no lineal y la decisión pierde hasta 0.33~K de error equivalente. N2 relinealiza a lo
largo de la trayectoria y no tiene ese problema.

\takeaway{\textbf{Respuesta a Q7.} Sí. N1 y N2 resuelven su QP convexo al óptimo global, y en
3800 consultas de lazo cerrado N3 queda a $\le2.5\cdot10^{-8}$~K y N2 a $\le6\cdot10^{-4}$~K de
error equivalente del mejor plan encontrado; un arranque en frío nunca encontró algo mejor. N1
se aleja (hasta 0.33~K) sólo con movimientos grandes del actuador. Es evidencia de optimalidad
local con varios arranques, no un certificado global: el certificado sigue limitado a las dos
consultas de la sección~\ref{sec:certificacion}.}
